% ============================================================================
% Guided MATLAB coding tutorial: indirect + direct solvers for the
% minimum-time low-thrust GTO -> tulip transfer (CR3BP).
% Compile: /Library/TeX/texbin/pdflatex (triple-pass for the TOC).
% ============================================================================
\documentclass[11pt,letterpaper]{article}
\usepackage[margin=1in]{geometry}
\usepackage{amsmath,amssymb,bm}
\usepackage{booktabs}
\usepackage{enumitem}
\usepackage{xcolor}
\usepackage[colorlinks=true,linkcolor=blue!70!black,urlcolor=blue!70!black]{hyperref}
\usepackage{fancyhdr}
\usepackage{titlesec}
\usepackage{mathtools}
\usepackage{tcolorbox}

\pagestyle{fancy}
\fancyhf{}
\rhead{GTO $\to$ Tulip Min-Time Solvers --- Guided Exercises}
\lhead{}
\rfoot{Page \thepage}
\setlength{\headheight}{14pt}

\setlength{\parskip}{5pt}
\setlength{\parindent}{0pt}

\titlespacing*{\section}{0pt}{14pt}{6pt}
\titlespacing*{\subsection}{0pt}{10pt}{4pt}
\titlespacing*{\subsubsection}{0pt}{8pt}{3pt}

% --- convenience macros ---
\newcommand{\R}{\mathbb{R}}
\newcommand{\T}{^{\mathsf{T}}}
\newcommand{\half}{\tfrac12}
\newcommand{\rvec}{\mathbf{r}}
\newcommand{\vvec}{\mathbf{v}}
\newcommand{\lamr}{\bm{\lambda}_r}
\newcommand{\lamv}{\bm{\lambda}_v}
\newcommand{\lamm}{\lambda_m}

% --- Exercise: what the learner writes (blue) ---
\newtcolorbox{exercise}[1]{
  colback=blue!5!white, colframe=blue!50!black, fonttitle=\bfseries,
  title={Exercise: #1}, sharp corners, boxrule=0.8pt,
  left=6pt, right=6pt, top=4pt, bottom=4pt}

% --- Hint: partial code / gotcha nudge (yellow) ---
\newtcolorbox{hint}{
  colback=yellow!5!white, colframe=yellow!60!black, fonttitle=\bfseries,
  title={Hint}, sharp corners, boxrule=0.5pt,
  left=6pt, right=6pt, top=4pt, bottom=4pt}

% --- Checkpoint: a runnable test with VERIFIED expected output (green) ---
\newtcolorbox{checkpoint}[1]{
  colback=green!5!white, colframe=green!50!black, fonttitle=\bfseries,
  title={Checkpoint: #1}, sharp corners, boxrule=0.8pt,
  left=6pt, right=6pt, top=4pt, bottom=4pt}

\title{\textbf{Building the GTO $\to$ Tulip Minimum-Time Solvers}\\[8pt]
\large Guided Coding Exercises\\[4pt]
\normalsize An indirect (PMP shooting) solver and a direct (collocation NLP)
solver for the low-thrust cislunar transfer behind pumpkynPie's
\texttt{lowThrust\_GTO\_Tulip.m}}
\author{}
\date{July 2026}

\begin{document}
\maketitle
\thispagestyle{fancy}

\tableofcontents
\newpage

% ==========================================================================
\section{Introduction}

You will build, from scratch, both classical solution machines for a real
trajectory-optimization problem: the 27.9-day minimum-time low-thrust spiral
from GTO to a 7-petal southern tulip orbit in the Earth--Moon CR3BP (15\,kg,
25\,mN, $I_{sp}=2100$\,s). The \textbf{indirect} solver integrates
Pontryagin's necessary conditions and shoots on the unknown initial
costates; the \textbf{direct} solver transcribes the trajectory onto a mesh
and hands up to $10^5$ variables to \texttt{fmincon}. Both must reproduce
the pumpkyn reference answer:
\[
  t_f = 6.290694\ \text{ND} = 27.8845\ \text{days},\qquad
  m_{\text{prop}} = 2.9247\ \text{kg},\qquad
  \Delta V = 4.4665\ \text{km/s}.
\]
The two ``aha''s you should reach: (i) on a 40-revolution spiral,
\emph{sensitivity is the whole game} --- it dictates complex-step
derivatives on the indirect side and density-matched meshes on the direct
side; (ii) the two methods are one problem wearing two costumes, and they check
each other to mesh accuracy ($\sim 3\times 10^{-4}$ relative in $t_f$ at
$N = 12{,}000$, improving under refinement).

\subsection{Prerequisites}

\begin{itemize}[leftmargin=1.6em,nosep]
  \item The companion theory note \texttt{gto\_tulip\_mintime\_theory.pdf}
  (this folder): the OCP, the PMP conditions, the transcription. Read it
  first; this document assumes its notation and derivations.
  \item \texttt{orbit\_transfer\_exercises.pdf}
  (\texttt{optimal\_control/orbit\_transfer/}): the two-body min-energy
  version of exactly this indirect/direct pairing. This tutorial is its
  CR3BP min-time sequel --- harder in every dimension that matters.
  \item The CoV/PMP notes \texttt{cov\_pmp\_orbit\_transfer.pdf}
  (\texttt{myLatex/notes/}) for where the costate ODE and transversality
  conditions come from.
  \item A working \texttt{pumpkyn} clone at
  \texttt{proj7/external/pumpkyn} (used for problem setup and as the
  reference implementation to test against).
\end{itemize}

\subsection{What you will build}

Two folders; the min-time build (Phases A--F), then the min-fuel
extension (Phases G--H):

\begin{center}
\small
\begin{tabular}{llp{6.1cm}}
\toprule
folder & file & purpose \\
\midrule
\texttt{lowThrust\_GTO\_tulip/}
 & \texttt{lt\_pmp\_eom.m} & 14-state augmented PMP dynamics (state +
   costate + optimal control law) \\
 & \texttt{shoot\_residual\_tf.m} & terminal-condition residual + Jacobian
   (complex step + analytic $t_f$ column) \\
 & \texttt{solve\_tfmin\_indirect.m} & \texttt{fsolve} shooting driver \\
 & \texttt{run\_gto\_tulip\_indirect.m} & endpoints via pumpkyn, solve,
   $\Delta V$ accounting, two-panel figure \\
\midrule
\texttt{NLP\_lowThrust\_GTO\_tulip/}
 & \texttt{lt\_dynamics.m} & 7-state control-explicit dynamics + analytic
   Jacobians $A, B$ \\
 & \texttt{unpack\_z.m} & decision-vector bookkeeping \\
 & \texttt{nlp\_constraints.m} & trapezoidal defects + unit-sphere control
   equalities, sparse gradient \\
 & \texttt{build\_guess.m} & density-matched mesh + warm start \\
 & \texttt{solve\_tfmin\_nlp.m} & \texttt{fmincon} interior-point driver \\
 & \texttt{NLP\_lowThrust\_GTO\_Tulip.m} & direct-method twin of the
   pumpkynPie demo \\
\midrule
min-fuel (G--H)
 & \texttt{lt\_pmp\_eom\_minfuel.m} & smoothed min-fuel PMP dynamics \\
 & \texttt{shoot\_residual\_minfuel.m} & fixed-$t_f$ 7-condition residual \\
 & \texttt{solve\_minfuel\_indirect.m} & LM + smoothing continuation \\
 & \texttt{lt\_dynamics\_throttle.m} & 4-control dynamics (throttle) \\
 & \texttt{nlp\_constraints\_minfuel.m} & fixed-$t_f$ defects + cone \\
 & \texttt{solve\_minfuel\_nlp.m} & max-$m_f$ NLP, fixed $t_f$ \\
 & \texttt{NLP\_lowThrust\_GTO\_Tulip\_minfuel.m} & leg-replan driver
   (direct $\to$ seed $\to$ polish attempt) \\
 & \texttt{costate\_seed\_from\_nlp.m} & LSQ costate reconstruction \\
\bottomrule
\end{tabular}
\end{center}

\subsection{Workspace split}

Build your versions in \texttt{lowThrust\_GTO\_tulip/mytry/} (create an
\texttt{NLP/} subfolder there for Phase D onward, or a second
\texttt{mytry/} inside \texttt{NLP\_lowThrust\_GTO\_tulip/}). The parent
folders hold the reference solutions --- consult them only \emph{after}
attempting each exercise. Run \texttt{setup\_paths} (in either folder) once
per session to put pumpkyn on the path.

\subsection{Notation recap}

ND units: length $l^\star = 389{,}703$\,km, time $t^\star = 382{,}981$\,s,
mass $m_0 = 15$\,kg; $\mu^\star = 0.012150585609624$. State
$x = (\rvec,\vvec,m) \in \R^7$; costates
$\lambda = (\lamr,\lamv,\lamm)$. ND thrust acceleration
$T_{\max} = 0.627292$, exhaust velocity $c = 20.238741$. Hamiltonian
$H = 1 + \lambda\T f$; primer direction $-\lamv/\|\lamv\|$; switching
function $S = -\|\lamv\| c/m - \lamm$. NLP decision vector
$Z = [X(:); W(:); t_f]$ on mesh fractions $\sigma_k$. The demo endpoints
(GTO departure, tulip arrival) are 6-vectors $rv_0, rv_f$ built by pumpkyn
calls in the drivers.

\begin{hint}
Complex-step house rules apply to \emph{every} dynamics file you write
here, same as \texttt{mlabTools/track\_lib}: magnitudes as
\texttt{sqrt(sum(x.\^{}2))} never \texttt{norm}, transposes as \texttt{.'}
never \texttt{'}, no \texttt{abs}, no \texttt{dot}. In Phase B your
Jacobian will differentiate \emph{through the integrator} with a
$10^{-20}$ imaginary perturbation --- one \texttt{norm} anywhere in the
call chain silently zeroes the derivative.
\end{hint}

% ==========================================================================
\section{Phase A: the augmented PMP dynamics}

The indirect method needs one workhorse: the 14-dimensional coupled
state--costate flow with the optimal control law evaluated pointwise from
the costates (theory note \S3). Everything else is bookkeeping around it.

\begin{exercise}{\texttt{lt\_pmp\_eom.m}}
Create
\begin{verbatim}
function [yDot, Ht, S, aThrust] = lt_pmp_eom(~, y, Tmax, c, muStar)
\end{verbatim}
with \texttt{y = [r(3); v(3); m; lambda\_r(3); lambda\_v(3); lambda\_m]}.
Steps:
\begin{enumerate}[leftmargin=1.5em,nosep]
\item Distances to the primaries: $\mathbf d = \rvec + (\mu^\star,0,0)$,
      $\mathbf s = \rvec - (1{-}\mu^\star,0,0)$; build $g(\rvec)$
      (centrifugal + two gravity terms) and the Coriolis vector
      $h(\vvec) = (2\dot y, -2\dot x, 0)$.
\item The two Jacobians you already know from the theory note: the
      symmetric $G = \partial g/\partial\rvec$ assembled in matrix form
      ($\mathrm{diag}(1,1,0)$ minus the two
      $I/\rho^3 - 3\bm\rho\bm\rho\T/\rho^5$ gravity-gradient blocks) and
      the constant Coriolis skew $H_c$.
\item Control law: $\|\lamv\|$, primer direction
      $\bm\alpha = -\lamv/\|\lamv\|$, switching function
      $S = -\|\lamv\| c/m - \lamm$, throttle $u = 1$ unless $S > 0$.
\item Dynamics: $\dot\rvec = \vvec$;
      $\dot\vvec = g + h + u\,(T_{\max}/m)\,\bm\alpha$;
      $\dot m = -u\,T_{\max}/c$; costates
      $\dot\lamr = -G\T\lamv$, $\dot\lamv = -\lamr - H_c\T\lamv$,
      $\dot\lamm = -\|\lamv\|\,u\,T_{\max}/m^2$.
\item \texttt{Ht} = $1 + \lamr\T\dot\rvec + \lamv\T\dot\vvec + \lamm\dot m$
      (guard with \texttt{nargout}); the fourth output is simply the
      thrust-acceleration vector,
      \texttt{aThrust} $= u\,(T_{\max}/m)\,\bm\alpha$, handy for plotting.
\end{enumerate}
Keep it complex-safe (see the Introduction hint). The reference
implementation you are re-deriving is \texttt{pumpkyn.cr3bp.tfMinEoM} ---
do \emph{not} peek at it until Checkpoint A passes; it also carries the
$14\times14$ analytic STM you are deliberately \emph{not} building
(Phase B replaces it with ten lines of complex step).
\end{exercise}

\begin{checkpoint}{A --- EoM agreement with pumpkyn}
\begin{verbatim}
setup_paths();
muStar = 0.012150585609624; tStar = 382981.289129055;
lStar  = 389703.264829278;
Tmax = (0.025/15)*tStar^2/(lStar*1000);   % 0.627292100855
c    = (2100/tStar)*(9.80665*tStar^2/(1000*lStar));
rv0  = [0.00349629072295 -0.00729625826008 0 ...
        4.19147893803 8.98865558978 0];
lam0 = [190.476047002 -79.7060409 -0.4298691 ...
        0.3011592775 0.5866700046 -0.0071173489 4.329378839]';
y  = [rv0(:); 1; lam0];
[yDotMine, HtMine, S] = lt_pmp_eom(0, y, Tmax, c, muStar);
yDotRef = pumpkyn.cr3bp.tfMinEoM(0, y, Tmax, c, muStar);
max(abs(yDotMine - yDotRef(1:14)))
S
y2 = y; y2(1:3) = [0.5;0.1;0.01]; y2(4:6) = [0.2;0.9;-0.01]; y2(7) = 0.9;
[yd1] = lt_pmp_eom(0, y2, Tmax, c, muStar);
yd2 = pumpkyn.cr3bp.tfMinEoM(0, y2, Tmax, c, muStar);
max(abs(yd1 - yd2(1:14)))
\end{verbatim}
Verified expected output:
\begin{verbatim}
max|yDot diff| (GTO departure state) ~ 3e-11
S (departure)                        = -17.6767
max|yDot diff| (mid-transfer state)  ~ 4.4e-16
\end{verbatim}
The departure-state agreement is $10^{-11}$ \emph{absolute}, not
$10^{-16}$: at 350\,km perigee altitude the gravity-gradient entries reach
$\sim 10^{9}$ and your matrix-form $G$ rounds differently than pumpkyn's
expanded scalar polynomials --- relative agreement is still machine-level.
If instead you see $O(1)$ disagreement in rows 8--13: check the transpose
and sign in $\dot\lamr = -G\T\lamv$ and the coupling term $-\lamr$ in
$\dot\lamv$ (the classic ``dropped kinetic coupling'' trap). Row 14 wrong
$\Rightarrow$ revisit $\dot\lamm$'s $m^2$.
\end{checkpoint}

% ==========================================================================
\section{Phase B: the shooting residual and its Jacobian}

The TPBVP (theory note \S3): find $z = [\lambda(0); t_f] \in \R^8$ so that
integrating the augmented flow from the (fixed) initial state produces
\[
  R(z) = \bigl[\rvec(t_f)-\rvec_f;\ \vvec(t_f)-\vvec_f;\ \lamm(t_f);\
  H(t_f)\bigr] = 0 .
\]

\begin{exercise}{\texttt{shoot\_residual\_tf.m}}
Create
\begin{verbatim}
function [R, J] = shoot_residual_tf(z, rv0, rvf, Tmax, c, muStar)
\end{verbatim}
\begin{enumerate}[leftmargin=1.5em,nosep]
\item Residual: integrate \texttt{lt\_pmp\_eom} over $[0, z_8]$ from
      \texttt{[rv0(:); 1; z(1:7)]} with \texttt{ode113} at
      \texttt{RelTol 1e-12 / AbsTol 1e-14}, evaluate the four blocks above
      (call the EoM once more at $y_f$ for $H(t_f)$).
\item Jacobian columns 1--7 by \textbf{complex step}: perturb each costate
      by $ih\,\max(1,|z_k|)$ with $h = 10^{-20}$, re-integrate, take
      $\mathrm{imag}(R)/(h\cdot\text{scale})$. (\texttt{ode113} propagates
      complex states without complaint as long as your EoM is
      complex-safe.)
\item Column 8 analytically: the terminal conditions ride the flow, so
      $\partial R/\partial t_f =
      [\dot\rvec;\dot\vvec;\dot\lamm;\dot H](t_f)$ --- and $\dot H = 0$
      exactly (autonomous canonical flow). One EoM call, no integration.
\end{enumerate}
\end{exercise}

\begin{hint}
Why not finite differences? Try it and watch: this arc crosses perigee
$\sim$40 times and amplifies initial perturbations by $\sim 10^6$, so the
residual's integration-noise floor sits far above any workable FD step.
With central FD, \texttt{fsolve}'s dogleg rejects every trial step ---
first-order optimality frozen at 487, ``locally singular'', zero progress
from the exact same guess that converges in five iterations with the
complex-step Jacobian. The theory note's Key Fact box has the full story.
Also: make the bang-bang branch explicitly complex-step-safe --- write
\texttt{if real(S) > 0} rather than relying on MATLAB's (real-part, but
easy-to-misread) complex comparison semantics.
\end{hint}

\begin{checkpoint}{B --- residual at the pumpkyn solution}
Rebuild the endpoints in full precision (copy the endpoint-construction
block from the reference \texttt{run\_gto\_tulip\_indirect.m}, or your own
Phase-C driver), then evaluate at pumpkyn's converged solution:
\begin{verbatim}
zRef = [190.476047002 -79.7060401966 -0.429868646477 ...
        0.301159276271 0.58667000078 -0.00711734996094 ...
        4.3293789063 6.2906939607]';
[R, J] = shoot_residual_tf(zRef, rv0, rvf, Tmax, c, muStar);
norm(R), R(7), R(8)
J(:,8)'   % the analytic tf column
\end{verbatim}
Verified expected output (endpoints built in full precision --- pasting
truncated 12-digit endpoint values roughly doubles this floor):
\begin{verbatim}
||R||  ~ 4e-5   (sensitivity-limited noise floor)
R(7)   ~ -3e-7,  R(8) ~ 5e-10
J(:,8) = [-0.0016 1.155 -0.0027 -25.6 -0.81 -42.8 -0.17 0]
\end{verbatim}
Read that last column like an astrodynamicist: rows 1--3 are the arrival
velocity (it IS $\dot\rvec(t_f) = \vvec_f$), rows 4--6 the $\sim$43-ND
near-Moon gravity at the tulip point, row 8 exactly 0 ($\dot H = 0$).
One full \texttt{[R, J]} evaluation costs 8 integrations ($\sim$2\,s).
If \texttt{J} columns 1--7 come back all zero: something in your call
chain destroyed the imaginary part --- hunt the \texttt{norm}/\texttt{'}
/\texttt{abs}.
\end{checkpoint}

% ==========================================================================
\section{Phase C: the indirect solve}

\begin{exercise}{\texttt{solve\_tfmin\_indirect.m} and
\texttt{run\_gto\_tulip\_indirect.m}}
\begin{enumerate}[leftmargin=1.5em,nosep]
\item \texttt{solve\_tfmin\_indirect}: wrap \texttt{fsolve}
      (\texttt{trust-region-dogleg}, \texttt{SpecifyObjectiveGradient
      true}, \texttt{StepTolerance 1e-12}, $\sim$400 function-evaluation
      budget) around your residual. Return the solution, $\|R\|$, and the
      exit flag.
\item \texttt{run\_gto\_tulip\_indirect}: reproduce the pumpkynPie demo
      end-to-end --- build the GTO state
      (\texttt{orb2eci}$\to$\texttt{fromPCI}), the tulip arrival
      (\texttt{getTulip}$\to$\texttt{prop}, arrival index = max $\dot y$),
      nondimensionalize the spacecraft, solve from the demo's published
      \texttt{lambda0\_tf} guess, re-propagate, and report $t_f$,
      propellant $m_0(1 - m(t_f))$, and
      $\Delta V = c\,\ln(1/m(t_f))\, l^\star/t^\star$. Add the two-panel
      (top-down / edge-on) trajectory figure with one coast period of the
      tulip appended.
\end{enumerate}
\end{exercise}

\begin{checkpoint}{C1 --- the transfer, solved}
\texttt{out = run\_gto\_tulip\_indirect(true);} Verified expected output
($\sim$25\,s total):
\begin{verbatim}
tf         = 6.290694 ND = 27.8845 days
||R||      ~ 1e-6
propellant = 2.9247 kg of 15 kg,  dV = 4.4665 km/s
\end{verbatim}
$t_f$ agrees with \texttt{pumpkyn.cr3bp.tfMin} (6.2906939607) to
\emph{eight} significant figures. Expect exit flag $-3$ (``locally
singular'') \emph{with} a tiny residual: \texttt{fsolve} has hit the
sensitivity-limited noise floor where no step improves $\|R\|$ --- the
equation is solved even though the flag sulks. The reference toolbox run
stalls the same way at $\|f\|^2 = 3\times10^{-17}$.
\end{checkpoint}

\begin{checkpoint}{C2 --- the basin is tiny (run it, feel it)}
Perturb the converged solution by 0.1\% on the costates and $+0.05$ on
$t_f$ and re-solve. Verified result: \texttt{fsolve} exhausts its budget
at $\|R\| = 0.74$ with $t_f$ wandering past 7.1 --- \emph{divergence from
a 0.1\% perturbation}. That is single shooting on a $10^6$-sensitivity
problem: the demo ships converged costates because nothing less converges.
(Fixes, in increasing effort: continuation in $T_{\max}$ or endpoint,
multiple shooting, or Phase D's direct method --- which is exactly why it
exists.)
\end{checkpoint}

% ==========================================================================
\section{Phase D: control-explicit dynamics for the NLP}

Move to \texttt{NLP\_lowThrust\_GTO\_tulip/}. The direct method needs the
dynamics as $f(x, w)$ with the control an \emph{argument} (no costates
anywhere), plus the Jacobians $A = \partial f/\partial x$,
$B = \partial f/\partial w$ for the defect gradient. One modeling decision
is already made for you (theory note \S4): the throttle is \emph{fixed} at
1 --- Phase A's switching function stays in $[-46.8, -2.5]$, so the
min-time arc never coasts, and carrying a throttle variable whose warm
start sits exactly on its bound derails interior-point methods (we
verified: \texttt{fmincon} shoves the throttle off the wall, coasts
0.38\,ND, and stalls 2\% suboptimal). Control = unit direction vector
$\mathbf w$, enforced by $\mathbf w\T\mathbf w = 1$ equalities.

\begin{exercise}{\texttt{lt\_dynamics.m} and \texttt{unpack\_z.m}}
\begin{verbatim}
function [F, A, B] = lt_dynamics(X, W, Tmax, c, muStar)   % X 7xM, W 3xM
function [X, W, tf] = unpack_z(Z, N)
\end{verbatim}
\begin{enumerate}[leftmargin=1.5em,nosep]
\item Vectorize $f$ across nodes (columns): reuse Phase A's $g, h$ math;
      thrust acceleration $T_{\max}\mathbf w/m$; mass flow constant
      $-T_{\max}/c$.
\item Per-node Jacobians in a loop: $A$ gets $I$ (position rows), $G$,
      $H_c$, and the mass column $-T_{\max}\mathbf w/m^2$; $B$ gets
      $(T_{\max}/m) I$ in the velocity rows. Return as
      \texttt{[7x7xM]}, \texttt{[7x3xM]}.
\item \texttt{unpack\_z}: $Z = [X(:); W(:); t_f]$,
      $n_Z = 10(N{+}1)+1$.
\end{enumerate}
\end{exercise}

\begin{checkpoint}{D --- consistency with Phase A}
With \texttt{y} the converged departure state from Checkpoint A and
$\mathbf w$ the primer direction computed from its costates:
\begin{verbatim}
lamv = y(11:13); w = -lamv/sqrt(sum(lamv.^2));
F = lt_dynamics(y(1:7), w, Tmax, c, muStar);
yDot = lt_pmp_eom(0, y, Tmax, c, muStar);
max(abs(F - yDot(1:7)))
\end{verbatim}
Verified expected output: exactly \texttt{0} (same arithmetic, same
order; anything visible means your thrust or mass-flow term differs from
Phase A's).
\end{checkpoint}

% ==========================================================================
\section{Phase E: transcription --- defects, sphere, and the mesh that matters}

\begin{exercise}{\texttt{nlp\_constraints.m} and \texttt{build\_guess.m}}
\begin{enumerate}[leftmargin=1.5em,nosep]
\item \texttt{nlp\_constraints(Z, sigma, Tmax, c, muStar)}: node times are
      $t_k = t_f\,\sigma_k$ with \emph{fixed} fractions $\sigma$
      (so free $t_f$ scales all segment widths $h_k = t_f\,\Delta\sigma_k$
      smoothly). Equalities: $7N$ trapezoidal defects
      $x_{k+1}-x_k-\frac{h_k}{2}(f_k+f_{k+1})$ and $N{+}1$ sphere
      constraints $\mathbf w_k\T\mathbf w_k - 1$. Assemble the sparse
      gradient from the per-segment blocks
      ($-I-\frac{h_k}{2}A_k$, $I-\frac{h_k}{2}A_{k+1}$,
      $-\frac{h_k}{2}B_{(\cdot)}$, and the $t_f$ column
      $-\frac{\Delta\sigma_k}{2}(f_k+f_{k+1})$) as triplet lists;
      remember \texttt{fmincon} wants the \emph{transpose} ($n_Z \times
      n_{ceq}$).
\item \texttt{build\_guess('indirect', N, ...)}: solve the indirect
      problem via \texttt{pumpkyn.cr3bp.tfMin}, propagate with
      \texttt{tfMinProp}, and --- the load-bearing trick --- build
      $\sigma$ as \textbf{quantiles of the integrator's own adaptive time
      grid} (equal integrator-steps per segment), so mesh density follows
      the dynamics. Interpolate states and primer directions onto
      $\sigma t_f$ (\texttt{pchip}; \texttt{unique} the ode abscissae
      first).
\end{enumerate}
\end{exercise}

\begin{hint}
Validate the sparse gradient the cheap way before trusting it at
$N=3000$ --- exactly this recipe:
\begin{verbatim}
rng(7); N = 5; nN = N+1; sigma = sort([0; rand(N-1,1); 1]);
Z = zeros(10*nN+1, 1);
for k = 1:nN
  Z((k-1)*7+(1:7)) = [0.4+0.1*randn(3,1); 0.5*randn(3,1); 0.8+0.1*rand];
  w = randn(3,1);  Z(7*nN+(k-1)*3+(1:3)) = w/norm(w)*0.95;
end
Z(end) = 5.5;   % then central FD on ceq with step 1e-7, all 61 columns
\end{verbatim}
(the random states sit near $r \approx 0.4$, safely away from both
primaries). Verified: $n_{ceq} = 41$, $n_Z = 61$, max relative Jacobian
error \texttt{1.6e-9}, 323 nonzeros. This is the checkpoint that catches
transposed blocks, wrong $t_f$ columns, and off-by-one node indexing ---
at $N=5$ you can print the full matrices and \emph{look}.
\end{hint}

\begin{checkpoint}{E --- the mesh is the message}
Warm-start defect $\|\cdot\|_\infty$ at the interpolated indirect
solution, verified:
\begin{verbatim}
                     uniform mesh      density-matched mesh
N = 1500                  --                 4.9e-3
N = 3000                0.90                 6.4e-4
N = 6000                  --                 8.1e-5
N = 12000                 --                 1.6e-5
\end{verbatim}
A \emph{uniform} 3000-segment mesh leaves defects of order one --- the
mesh step (13.4\,min) barely samples the 3.5\,$\times 10^{-3}$-ND perigee
passes, and \texttt{fmincon} inherits garbage. The density-matched mesh is
1400$\times$ better at identical cost, and the table drops $\sim 8\times$
per mesh doubling --- the \emph{local} trapezoidal defect rate $O(h^3)$
(global trajectory accuracy is $O(h^2)$). If your density-matched numbers
are much worse: your $\sigma$ probably has equal \emph{time} spacing, not
equal \emph{integrator-step} spacing.
\end{checkpoint}

% ==========================================================================
\section{Phase F: the NLP solve}

\begin{exercise}{\texttt{solve\_tfmin\_nlp.m} and
\texttt{NLP\_lowThrust\_GTO\_Tulip.m}}
\begin{enumerate}[leftmargin=1.5em,nosep]
\item \texttt{solve\_tfmin\_nlp}: objective $J(Z) = Z_{\text{end}}$ with a
      sparse one-hot gradient; pin $x_1$ and the first six components of
      $x_{N+1}$ by \texttt{lb = ub}; loose safety bounds everywhere else
      (states $\pm 3/\pm 12$, mass $[0.3, 1]$, $\mathbf w \in
      [-1.5,1.5]^3$, $t_f \in [2, 15]$) --- \emph{every} bound strictly
      inactive at the warm start. \texttt{fmincon} interior-point with
      \texttt{SpecifyObjectiveGradient}, \texttt{SpecifyConstraintGradient},
      \texttt{HessianApproximation 'lbfgs'}, and
      \texttt{InitBarrierParam 1e-6} (the default 0.1 barrier yanks a
      warm-started iterate far off the solution before pulling it back ---
      we watched it drive $t_f$ to 4.0 at feasibility 1.1).
\item \texttt{NLP\_lowThrust\_GTO\_Tulip(N, guessMode, makePlot)}: same
      endpoint setup as Phase C's driver, then guess $\to$ solve $\to$
      report $t_f$/propellant/$\Delta V$, max defect, and the node-wise
      deviation from the indirect arc, plus the demo-style two-panel
      figure. Do \emph{not} ``validate'' by flying the control open-loop
      through \texttt{ode113}: on this problem replay diverges ($10^5$\,km)
      for \emph{any} control including the true optimum --- interpolation
      error $\times\ 10^6$ amplification. Node-wise trajectory comparison
      is the honest check.
\end{enumerate}
\end{exercise}

\begin{checkpoint}{F --- the direct answer}
\texttt{out = NLP\_lowThrust\_GTO\_Tulip(12000, 'indirect', true);}
Verified expected output ($\sim$1\,min; fmincon exit flag 2, step-size
stall at the solution --- same species as Phase C's flag $-3$):
\begin{verbatim}
tf         = 6.288574 ND = 27.8751 days   (indirect: 6.290694 / 27.8845)
max defect = 1.7e-15
propellant = 2.9237 kg,  dV = 4.4648 km/s (indirect: 2.9247 / 4.4665)
max node deviation from indirect arc: 1.0e+03 km, 0.023 km/s
\end{verbatim}
And the part worth staring at --- the mesh-refinement convergence of the
\emph{discrete} optimum toward the continuous one (each $\sim$0.5--1\,min):
\begin{verbatim}
N =  3000:  tf ~ 6.2658    (error -2.5e-2, node dev 11,000 km)
N =  6000:  tf = 6.281802   (error -8.9e-3, node dev  3,970 km)
N = 12000:  tf = 6.288574   (error -2.1e-3, node dev  1,010 km)
\end{verbatim}
(Quote the $N{=}3000$ value loosely: at coarse-mesh accuracy the stall
point scatters by $\sim 10^{-3}$ across numerically equivalent code
variants --- the valley is that flat. It sharpens as $N$ grows: the finer
meshes reproduce to 7 digits.)
Defects sit at machine zero at EVERY N --- the discrete dynamics are
solved exactly --- yet $t_f$ is \emph{below} the true optimum and climbs
back as the mesh refines. The optimizer is (legitimately!) exploiting
trapezoidal discretization error near perigee: coarse-mesh physics is
slightly wrong physics, and a minimizer will always find the flaw. This is
why direct-method answers are quoted \emph{with} a refinement study, never
from one mesh.
If \texttt{fmincon} instead wanders off (objective plunging while
feasibility grows): you are seeing the warm-start/barrier failure modes
this phase's settings exist to prevent --- check \texttt{InitBarrierParam}
and that no bound is active at $Z_0$.
\end{checkpoint}

% ==========================================================================
\section{Phase G: min-fuel --- the machinery, and the trap}

Everything so far never used the switching function --- min-time burns
always. Min-fuel (theory note \S6) is posed on the ARRIVAL LEG as a
mid-transfer replan: start from the min-time arc's state at
$\tau = 4.0$ (with its remaining mass $m_0 = 0.876$), leg time fixed at
$1.3\times$ the leg's minimum ($t_f = 2.9779$ ND $= 13.20$ days), arrival
at the tulip point one ballistic coast downstream. The throttle law
becomes bang-bang on $S = 1 - \|\lamv\|c/m - \lamm$, smoothed for
numerics by $u_\varepsilon = (1 - \tanh(S/2\varepsilon))/2$.

\begin{exercise}{\texttt{lt\_pmp\_eom\_minfuel.m}, \texttt{shoot\_residual\_minfuel.m}, \texttt{solve\_minfuel\_indirect.m}}
\begin{enumerate}[leftmargin=1.5em,nosep]
\item Copy \texttt{lt\_pmp\_eom.m}; add \texttt{epsSmooth} as a trailing
      parameter; replace the switching function and throttle with the
      min-fuel forms above; in \texttt{Ht}, \emph{replace} the min-time
      running cost $1$ with $\tfrac{T_{\max}}{c}u$ (the bare $1$ belongs
      to min-time only). Costate ODEs unchanged (envelope theorem ---
      the smoothed $u$ is an \emph{interior} minimizer of the smoothed
      $H$).
\item The residual has SEVEN unknowns and seven conditions:
      $[\rvec(t_f)-\rvec_f;\ \vvec(t_f)-\vvec_f;\ \lamm(t_f)]$, with
      $t_f$ fixed and the initial mass a parameter (the leg starts at
      $m_0 < 1$) --- no $H(t_f)$ row, no $t_f$ column. Jacobian by
      complex step exactly as in Phase B.
\item \texttt{solve\_minfuel\_indirect(rv0, m0, rvf, tf, lamGuess,
      Tmax, c, muStar, epsSchedule)}: wrap \texttt{lsqnonlin}
      (Levenberg--Marquardt, gradient on) in an
      $\varepsilon$-continuation loop, re-seeding each solve from the
      previous converged costates. Use LM rather than \texttt{fsolve}'s
      trust-region-dogleg: measured on this problem, LM makes $4\times$
      the residual progress per iteration budget.
\end{enumerate}
\end{exercise}

\begin{checkpoint}{G --- the saturation trap (run it, learn it)}
This diagnostic runs on the FULL transfer (everything is already in
hand from Phase C: its endpoints, $m_0 = 1$, $t_f = 1.2 \times 6.2907$
ND, and its converged costates $z(1{:}7)$ as the seed). Seed the
shooting with the \emph{raw} min-time costates and watch it stall at
$\|R\| = 1.55$: with $\|\lamv\|c \approx 13$ the switching function sits
near $-17$, $\tanh$ saturates, $\partial u/\partial S \sim 10^{-8}$, and
the Jacobian is blind in the throttle channel. Verify the scale anchor:
the min-fuel ``$+1$'' fixes an absolute costate gauge
($S(0) = -16.7$ raw; rescaling by $1/\|\lamv(0)\|c$ gives
$S(0) = -0.32$). Rescaling alone still stalls ($\|R\| = 0.83$) --- the
min-fuel costate \emph{ratio} $\lamr/\lamv$ differs structurally from
min-time, which is why Phase H builds the seed from the direct solution
instead.
\end{checkpoint}

% ==========================================================================
\section{Phase H: min-fuel --- direct finds, indirect (tries to) certify}

The throttle returns as a decision variable --- the general 4-control
transcription the theory note kept alive is now the one you build.

\begin{exercise}{\texttt{lt\_dynamics\_throttle.m}, \texttt{nlp\_constraints\_minfuel.m}, \texttt{solve\_minfuel\_nlp.m}, \texttt{costate\_seed\_from\_nlp.m}, driver}
\begin{enumerate}[leftmargin=1.5em,nosep]
\item \texttt{lt\_dynamics\_throttle}: control vector $[\mathbf w;\,s]$
      (reserving $u$ for the PMP throttle),
      thrust $T_{\max}\mathbf w/m$, mass flow $-T_{\max}s/c$; $B$ is now
      $7\times4$ with $\partial\dot m/\partial s = -T_{\max}/c$.
\item \texttt{nlp\_constraints\_minfuel}: $t_f$ is a \emph{parameter} ---
      $Z = [X(:); U(:)]$, $n_Z = 11(N{+}1)$, widths
      $h_k = t_f\Delta\sigma_k$ constants, no $t_f$ column. Cone
      $\mathbf w\T\mathbf w - s^2 = 0$, gradient $[2\mathbf w; -2s]$.
\item \texttt{solve\_minfuel\_nlp}: objective $-m_{N+1}$ (linear,
      one-hot gradient), $s \in [0,1]$, initial mass pinned at the leg's
      $m_0$.
\item \texttt{costate\_seed\_from\_nlp(tMesh, X, U, Tmax, c, muStar)}:
      the covector mapping done properly. Integrate the LINEAR costate
      ODE $\dot\lambda = -A(t)\T\lambda$ once along the NLP arc as a
      $7\times7$ transition matrix $\Psi(t)$ (49 stacked ODEs,
      \texttt{pchip}-interpolated $X, U$); stack the homogeneous rows
      $\lamv(t_k)\times\bm\alpha_k = \mathbf 0$ at $\sim$120 thinned
      burn nodes plus $\lamm(t_f) = 0$; take the smallest right
      singular vector; fix the sign (thrust must \emph{oppose}
      $\lamv$) and the scale ($S = 0$ at the first throttle switch).
\item The driver (\texttt{NLP\_lowThrust\_GTO\_Tulip\_minfuel.m})
      builds the BURN+COAST warm start: the min-time tail
      (throttle $\to 0.98$) followed by a ballistic coast (throttle
      $0.02$), targets the coast terminus, meshes density-matched on the
      concatenated arc, and sets $\mathbf w = s\,\bm\alpha$ so the cone
      holds exactly. Two constructions that fail instructively first: a
      time-stretched arc (defects $\sim 0.04$ everywhere, fmincon
      grinds) and pinning the original $\rvec_f$ (one $O(0.3)$-ND defect
      cliff; feasibility mode returns flag $-2$).
\end{enumerate}
\end{exercise}

\begin{checkpoint}{H1 --- the direct min-fuel answer}
Gradcheck first (same $N=5$ recipe as Phase E, throttle in $[0.3,0.9]$):
verified $n_{ceq} = 41$, $n_Z = 66$, max relative error \texttt{1.4e-9},
304 nonzeros; and \texttt{lt\_dynamics\_throttle} with
$u = [u_\varepsilon\bm\alpha;\,u_\varepsilon]$ matches
\texttt{lt\_pmp\_eom\_minfuel} rows 1--7 exactly (diff \texttt{0}). Then
\texttt{out = NLP\_lowThrust\_GTO\_Tulip\_minfuel(3000, 1.3, 4.0, true);}
Verified expected output:
\begin{verbatim}
warm-start max defect ~ 9.5e-05   (burn+coast: feasible by construction)
NLP: flag 2 (or 0), max defect ~ 2e-15
propellant = 1.0622 kg  (burn+coast reference: 1.0650)
burn fraction = 76.9%, single switch at t = 2.2907 = tfMinLeg
\end{verbatim}
The throttle plot is clean bang-bang: $s$ pinned at 1, one switch, then
pinned at 0 --- structure emerging from \emph{bound-active variables},
no switching function anywhere. The optimum trims 0.3\% from the
reference by throttle redistribution; with the phase-shifted target the
burn-then-coast profile is nearly optimal by construction.
\end{checkpoint}

\begin{checkpoint}{H2 --- the boss fight, quantified (expected: no convergence)}
The driver then reconstructs a shooting seed from the NLP solution with
your \texttt{costate\_seed\_from\_nlp} (exercise item 4) --- the
covector mapping done properly; by contrast, fmincon's own
multiplier estimates at an l-BFGS stall are $\sim$100$\times$ off in
scale (printed by the driver for comparison). Verified seed progression,
each strategy an order better, none inside the basin:
\begin{verbatim}
raw min-time costates (full)      ||R|| stalls at 1.55
rescaled min-time costates (full) ||R|| stalls at 0.83
switch-anchored min-time (leg)    ||R|| stalls at 0.33
LSQ costate reconstruction (leg)  ||R|| stalls at 0.14
\end{verbatim}
If your numbers land in this range, your implementation is
\emph{correct} --- the min-fuel shooting basin on a multi-revolution
CR3BP leg is simply smaller than the best seed available at this level
of machinery. Getting inside it takes switch-time (multi-arc) shooting
with event detection, direct multiple shooting, or a solver with
converged multipliers --- see What Comes Next.
\end{checkpoint}

% ==========================================================================
\section{What comes next}

\begin{enumerate}[leftmargin=1.6em]
\item \textbf{Close the min-fuel boss fight}: switch-time (multi-arc)
      shooting --- make the switch instants explicit unknowns with $S=0$
      interior conditions and event-detected integration between them ---
      or direct multiple shooting. These are the standard escalations
      for exactly this failure mode; neither is verified here, and
      converting Checkpoint H2's $\|R\| = 0.14$ into a converged
      indirect solution is the open exercise. The impulsive
      limit and Lawden's conditions (P-OC6) follow.
\item \textbf{Hermite--Simpson} defects: fourth-order accuracy, midpoint
      controls; the warm-start defect table of Checkpoint E should drop
      by orders of magnitude at fixed $N$.
\item \textbf{Adaptive mesh refinement}: re-mesh from your own solution's
      defect distribution (Betts Ch.~4) instead of borrowing the
      integrator's step density.
\item \textbf{Continuation}: solve at $10\times$ thrust (easy, few revs),
      walk $T_{\max}$ down to 25\,mN re-warm-starting each step --- the
      standard cure for Checkpoint C2's tiny basin, no published costates
      needed.
\item \textbf{Covector mapping}: extract the KKT multipliers of the defect
      constraints from \texttt{fmincon}'s \texttt{lambda.eqnonlin} and
      compare against the indirect costates along the arc --- the two
      methods' bookkeeping meets in one plot.
\end{enumerate}

\end{document}
